\subsection{Identidades con sus hipótesis}
\textbf{Combinaciones y Catalan.} $\binom nk=\binom n{n-k}$,
$\binom nk=\binom{n-1}{k-1}+\binom{n-1}k$, $\sum_k\binom nk=2^n$.
Catalan $C_0=1$, $C_n=\frac1{n+1}\binom{2n}n$ cuenta secuencias de paréntesis balanceados con $n$ pares y caminos que no cruzan la diagonal. También
$C_n=\sum_{i=0}^{n-1}C_iC_{n-1-i}$ y $C_n=\frac{2(2n-1)}{n+1}C_{n-1}$.
Para enteros exactos divide sin perder divisibilidad; en modular usa inversos sólo si existen.

\textbf{Distribuciones.} Soluciones $x_1+\cdots+x_k=n$, $x_i\ge0$:
$\binom{n+k-1}{k-1}$ para $k\ge1$; con $x_i\ge1$, $\binom{n-1}{k-1}$ si $n\ge k$. Para cotas superiores aplica inclusión-exclusión; no basta cambiar el factorial.

\textbf{Divisores.} Si $n=\prod p_i^{e_i}$,
$\tau(n)=\prod(e_i+1)$ es la \emph{cantidad} de divisores y
$\sigma(n)=\prod\frac{p_i^{e_i+1}-1}{p_i-1}$ es su \emph{suma}.
$\varphi(n)=n\prod_{p\mid n}(1-1/p)$ cuenta enteros en $[1,n]$ coprimos con $n$.
No hay cota $\tau(n)=O(\log\log n)$; usa la tabla de máximos al dimensionar enumeraciones.

\textbf{Möbius.} Si $g(n)=\sum_{d\mid n}f(d)$, entonces
$f(n)=\sum_{d\mid n}\mu(d)g(n/d)$.
$\sum_{d\mid n}\varphi(d)=n$ y $\sum_{d\mid n}\mu(d)=[n=1]$.
$\gcd(a,b)\operatorname{lcm}(a,b)=|ab|$; calcula el lcm como $(a/\gcd(a,b))b$ con tipo ancho.

\textbf{Fibonacci.} Con $F_0=0,F_1=1$:
$F_{2n+1}=F_n^2+F_{n+1}^2$,
$F_{2n}=F_n(2F_{n+1}-F_n)$,
$\sum_{i=0}^nF_i=F_{n+2}-1$.
Las dos primeras permiten fast doubling en $O(\log n)$; aplica módulo en cada operación y normaliza la resta.

\textbf{Esperanza y probabilidad.} $E[\sum_iX_i]=\sum_iE[X_i]$ sin exigir independencia.
Una variable indicadora vale 1 cuando ocurre el evento, luego $E[X_i]=P(evento_i)$.
Para estados con transición $p_{uv}$ y costo por paso $c_u$,
$E_u=c_u+\sum_vp_{uv}E_v$; terminal $E=0$.
Un ciclo de dependencias requiere resolver ecuaciones o justificar un orden; en modular las probabilidades requieren denominadores invertibles.

\textbf{Juegos imparciales finitos.} $g(s)=\operatorname{mex}\{g(t):s\to t\}$; mex es el menor entero no negativo ausente.
En juego normal (pierde quien no mueve), suma de juegos independientes es perdedora si el XOR de sus Grundy es 0. Nim es el caso $g(\text{montón})=\text{tamaño}$.
Debe terminar y tener iguales movimientos para ambos jugadores; misère o ciclos requieren otro análisis.
